% SPDX-License-Identifier: GPL-3.0-only
% Licensed under GNU GPL version 3; https://www.gnu.org/licenses/gpl-3.0.html . See LICENSE at repository root.
\section{State updates and symbolic labels}\label{sec:updates}

\subsection{Reusable state transformations}
Updates and rooted state expressions are different syntactic categories:
\begin{align*}
 r,q&\bnf\skipu\mid (x:=t)\mid r;q
           \mid\paru{(x_1,t_1),\ldots,(x_n,t_n)},\\
 U,V&\bnf i\mid\upd{r}{U}.
\end{align*}
An update contains no world names. A state expression starts from a world
name and applies updates inside out. The source constructors are
\code{Skip}, \code{Put}, \code{Then}, \code{Parallel}, \code{Base}, and
\code{Apply}. The compatibility spelling \code{AssignUpdate x t U} means
$\upd{x:=t}{U}$.
\sources{\rocq{Syntax.v}{update}, \rocq{Syntax.v}{state_expr},
  \rocq{Syntax.v}{AssignUpdate}}

Write $r(s)$ for the transformation of $s$ by $r$. Its meaning is
\begin{align*}
 \skipu(s)&=s,\\
 (x:=t)(s)&=s[x\leftarrow\eval{t}{s}],\\
 (r;q)(s)&=q(r(s)).
\end{align*}
For a parallel list, all right-hand sides are evaluated in the original input
$s$, then committed in list order. With $s_0=s$, set
\[
 s_j=s_{j-1}[x_j\leftarrow\eval{t_j}{s}],\qquad
 \paru{(x_1,t_1),\ldots,(x_n,t_n)}(s)=s_n.
\]
The rightmost duplicate target wins. The empty parallel update is identity.
\sources{\rocq{Semantics.v}{apply_update}, \rocq{Semantics.v}{eval_parallel},
  \rocq{Semantics.v}{parallel_last_write}}

\begin{example}[Snapshot reads]
Over integers, suppose $s(x)=1$ and $s(y)=2$ at their observed sort.
The parallel update $\paru{(x,y),(y,x)}$ gives $x=2,y=1$.
The sequential update $(x:=y);(y:=x)$ gives $x=2,y=2$.
Both outcomes follow by evaluating the displayed update definitions.
\end{example}

\begin{proposition}[Update algebra and frame]
Sequential update application is associative. Assigning a variable its
current value is identity. An update preserves the complete value of every
variable outside its syntactic target list.
\end{proposition}
\sources{\rocq{Semantics.v}{apply_then_associative},
  \rocq{Semantics.v}{apply_put_current}, \rocq{Semantics.v}{apply_update_frame}}

These semantic equations do not assert syntactic equality of update trees.
The public conversion requests implement the following four selected laws
inside state expressions:
\begin{align*}
 \upd{\skipu}{U}&\equiv U,&
 \upd{r;q}{U}&\equiv\upd{q}{\upd{r}{U}},\\
 \upd{\paru{}}{U}&\equiv U,&
 \upd{\paru{(x,t)}}{U}&\equiv\upd{x:=t}{U}.
\end{align*}
Here $\equiv$ means equal denotation in every world. It is not an instruction
to expand a general parallel list into sequential assignments.
\sources{\rocq{Syntax.v}{update_law}, \rocq{Syntax.v}{rewrite_state},
  \rocq{Semantics.v}{rewrite_state_sound}}

\subsection{Worlds and labels}
\begin{definition}[World and label]
A world $w$ gives a state $w_{\mathrm{s}}(i)$ to each world name and a nonempty
finite trace $w_{\mathrm{t}}(k)$ to each trace name. Labels have the grammar
\[
 \ell\bnf\epsilon\mid U\mid k\mid\ell\jlabel\ell.
\]
The state expression $U$ in label position stands for a singleton state label;
the trace name $k$ stands for a trace label.
\end{definition}
\sources{\rocq{Semantics.v}{world}, \rocq{Syntax.v}{label}}

State expressions and labels evaluate by
\begin{align*}
 \eval{i}{w}&=w_{\mathrm{s}}(i),&
 \eval{\upd{r}{U}}{w}&=r(\eval{U}{w}),\\
 \eval{\epsilon}{w}&=[],&
 \eval{U\text{ (as label)}}{w}&=[\eval{U}{w}],\\
 \eval{k}{w}&=w_{\mathrm{t}}(k),&
 \eval{\ell\jlabel\sigma}{w}&=\eval{\ell}{w}\cat\eval{\sigma}{w}.
\end{align*}
\sources{\rocq{Semantics.v}{eval_state}, \rocq{Semantics.v}{eval_label}}

Append neither checks equality of adjacent states nor merges them. Thus
$U\jlabel U$ denotes $[s,s]$, and $U\jlabel V$ denotes $[s,t]$ even when
$s\ne t$. A whole expression $\upd{r;q}{U}$ denotes one state, not the
sequence of intermediate states. To retain those states, write separate
state labels joined by $\jlabel$.

\begin{proposition}[Label algebra and nonemptiness]
The empty label is a two-sided semantic unit, and label juxtaposition is
semantically associative. The Boolean check \code{label_nonemptyb} is true
exactly when the evaluated label is nonempty, in every world.
It is false on $\epsilon$, true on a state or trace name, and uses disjunction
on a join.
\end{proposition}
\sources{\rocq{Semantics.v}{label_left_unit}, \rocq{Semantics.v}{label_right_unit},
  \rocq{Semantics.v}{eval_rotate_label}, \rocq{Semantics.v}{label_nonemptyb_spec}}

Empty components may therefore occur inside an admissible complete label.
No default state is used to interpret a wholly empty assertion label.

\subsection{Environment substitution and explicit conversions}
World substitution replaces a selected \code{Base} name by a state expression
throughout labels, leaving formulas untouched. Trace substitution replaces a
selected trace name by a label. Its replacement must be nonempty. The matching
semantic operation changes one mapping in $w$.
\sources{\rocq{Syntax.v}{subst_world_label}, \rocq{Syntax.v}{subst_trace_label},
  \rocq{Semantics.v}{eval_world_label_substitution},
  \rocq{Semantics.v}{eval_trace_label_substitution}}

The caller can select a Boolean path through label joins: false goes left,
true goes right. At that path, rotation reassociates a join, and unit
conversion inserts or removes an empty component. For update conversion,
the path selects a state label and a natural depth selects its enclosing
\code{Apply} layer. Update conversion returns failure for an invalid path,
depth, or law. Rotation and unit functions instead leave a nonmatching
subtree unchanged. None of these operations performs a normalization search.
\sources{\rocq{Syntax.v}{rotate_label}, \rocq{Syntax.v}{unit_label},
  \rocq{Syntax.v}{rewrite_update_label},
  \rocq{Semantics.v}{rewrite_update_label_sound}}
